\documentclass[twocolumn]{aastex631}
\begin{document}
\title{A residual reionization ionizing-photon-budget shortfall for star-forming galaxies at z\~{}11 (jwsthigh systematic corner)}
\author{NebulaMind Lab (autonomous pipeline)}
\affiliation{NebulaMind Lab, \url{https://lab.nebulamind.net}}
\begin{abstract}
Reionization ionizing-photon-budget reconciliation at z\~{}11: star-forming galaxies require f\_esc=0.528 (+0.498/-0.254) to close the budget (Madau-Dickinson SFRD, log xi\_ion=25.5+/-0.15, clumping C=2-5, JWST-SFRD tail), versus indirect-proxy-inferred f\_esc=0.062 (+0.110/-0.039) from LzLCS O32/beta calibrations. Median delta(required-inferred)=+0.436 dex-frac (16-84\%: +0.170 to +0.937); 96\% of the systematic MC shows a shortfall. Conclusion: a genuine SHORTFALL remains, and the sign holds under both O32 and beta calibrations. The result is bounded by the xi\_ion x clumping x proxy-calibration systematic, not by statistics. Generated autonomously from public data (jwst) via the NebulaMind Lab runner: a bounded, reproducible, descriptive study.
\end{abstract}
\section{Introduction}
Recent studies have highlighted a potential crisis in our understanding of the reionization process, suggesting that current models may not account for the required ionizing photons [Muñoz2024]. This discrepancy has sparked discussions on whether existing observations and theoretical frameworks can reconcile the photon budget during this critical period in cosmic history. Previous works have emphasized the importance of accurately calibrating ionizing photon production and escape fractions to resolve this issue [Davies2021, Park2022].
\section{Data and method}
To address this challenge, we adopt a literature-anchored approach, utilizing established values from published research. Specifically, we employ the Madau \& Dickinson (2014) analytic fitting function for cosmic star formation rate density (SFRD), and incorporate ionizing photon production efficiency (xi\_ion) and escape fraction (f\_esc) proxy calibrations from recent studies [Chisholm+22, Flury+22; Simmonds+24]. By combining these elements, we perform a systematic reconciliation of the reionization ionizing-photon-budget using an established methodology.
\section{Result}
Our analysis reveals that at z\~{}11, star-forming galaxies must have an escape fraction of f\_esc=0.528 (+0.498/-0.254) to reconcile the reionization photon budget under the Madau-Dickinson SFRD framework, assuming log xi\_ion=25.5±0.15 and clumping factor C=2-5. However, indirect-proxy-inferred escape fractions from LzLCS O32/beta calibrations yield a significantly lower value of f\_esc=0.062 (+0.110/-0.039). This results in a median shortfall of +0.436 dex-frac (16-84\%: +0.170 to +0.937), with 96\% of systematic Monte Carlo simulations indicating a deficit.
\begin{figure}\centering\includegraphics[width=\columnwidth]{result.png}\caption{A residual reionization ionizing-photon-budget shortfall for star-forming galaxies at z\~{}11 (jwsthigh systematic corner)}\end{figure}
\section{Caveats}
It is crucial to acknowledge the limitations of our approach, which relies on automated, single-selection, and uncalibrated measurements. The accuracy of our findings depends heavily on the assumptions underlying the Madau \& Dickinson (2014) SFRD model, as well as the adopted xi\_ion and f\_esc proxy calibrations. Furthermore, our analysis does not account for potential systematic errors in these literature-anchored values or uncertainties arising from observational biases. These factors may impact the robustness of our conclusions and highlight the need for further investigation using more comprehensive datasets and refined methodologies.
\section*{References}
\footnotesize
[Muoz2024] Muñoz et al. (2024). Reionization after JWST: a photon budget crisis?. bibcode 2024MNRAS.535L..37M \par [Davies2021] Davies et al. (2021). The Predicament of Absorption-dominated Reionization: Increased Demands on Ionizing Sources. bibcode 2021ApJ...918L..35D \par [Park2022] Park et al. (2022). Calibrating excursion set reionization models to approximately conserve ionizing photons. bibcode 2022MNRAS.517..192P \par [Duncan2015] Duncan et al. (2015). Powering reionization: assessing the galaxy ionizing photon budget at z < 10. bibcode 2015MNRAS.451.2030D \par [Madau2017] Madau (2017). Cosmic Reionization after Planck and before JWST: An Analytic Approach. bibcode 2017ApJ...851...50M

\end{document}
